\lysh{21308_SOLUTION}{1}
\textbf{ΛΥΣΗ}

α) Η εξίσωση της έλλειψης είναι της μορφής $\frac{x^2}{\alpha^2} + \frac{y^2}{\beta^2} = 1$ με $\alpha = 5$ και $\beta = 4$.

Επομένως έχει εστίες της μορφής $E(\gamma, 0)$ και $E'(-\gamma, 0)$, με $\gamma > 0$, όπου $\beta = \sqrt{\alpha^2 - \gamma^2}$ ή $\beta^2 = \alpha^2 - \gamma^2$ ή $\gamma^2 = \alpha^2 - \beta^2$.

Αντικαθιστώντας $\alpha = 5$ και $\beta = 4$ έχουμε $\gamma^2 = 5^2 - 4^2$ ή $\gamma^2 = 25 - 16$ ή $\gamma^2 = 9$ ή $\gamma = 3$, εφόσον $\gamma > 0$.

Άρα οι εστίες της έλλειψης είναι οι $E(3,0)$ και $E'(-3,0)$.

Η απόσταση των εστιών είναι $2\gamma = 2 \cdot 3 = 6$.

β) Ο μικρός άξονας της έλλειψης έχει μήκος $2\beta = 2 \cdot 4 = 8$.

Ο μεγάλος άξονας της έλλειψης έχει μήκος $2\alpha = 2 \cdot 5 = 10$.

γ) Από τη θεωρία η εφαπτομένη ($\varepsilon$) της έλλειψης της μορφής $\frac{x^2}{\alpha^2} + \frac{y^2}{\beta^2} = 1$ στο σημείο της $(x_1, y_1)$ είναι
\[
\frac{x \cdot x_1}{\alpha^2} + \frac{y \cdot y_1}{\beta^2} = 1.
\]
Αντικαθιστώντας όπου $x_1$ και $y_1$ τις συντεταγμένες του σημείου B της έλλειψης και όπου $\alpha = 5$ και $\beta = 4$ έχουμε:
\[
\frac{0x}{25} + \frac{4y}{16} = 1 \text{ ή } y = 4, \text{ που είναι η εξίσωση της } (\varepsilon).
\]
% TODO: TikZ σχήμα
\end{lysh}